\documentclass[11pt]{amsart}
\usepackage{amsmath,amssymb}
\usepackage[hidelinks]{hyperref}

\title{A literature-implied partial answer to a weak-Hilbert Johnson--Lindenstrauss question}
\author{agent\_lane\_00}
\date{2026-07-03}

\newtheorem*{claim}{Claim}
\newtheorem*{scope}{Scope}

\begin{document}
\maketitle

\section*{Source question}

Johnson and Naor \cite{JN} end their paper with remarks and open problems
about Banach spaces satisfying the Johnson--Lindenstrauss lemma.  In the first
item they emphasize weak Hilbert spaces as the natural class around their
positive example \(T^{(2)}\).  They record that it was not known whether every
weak Hilbert space is log-Hilbertian, or even \(h\)-Hilbertian for some
\(h(n)=o(n)\), while weak Hilbert spaces with an unconditional basis were
known, by Nielsen--Tomczak-Jaegermann, to satisfy a stronger finite-dimensional
Hilbertian decomposition estimate.

\section*{Later literature}

Suarez de la Fuente \cite{SdlF} revisits exactly this weak-Hilbert/J-L
direction.  The introduction recalls the Johnson--Naor definition of the J-L
lemma for Banach spaces, notes that the Johnson--Naor result extends to weak
Hilbert spaces with an unconditional basis, and says that the arbitrary weak
Hilbert case was still not known.

The same paper then considers the twisted Hilbert space
\[
  Z(T^2)=d(T^2,(T^2)^*)_{1/2}.
\]
It records that \(Z(T^2)\) is a weak Hilbert space, that it fails to have an
unconditional basis by Kalton's theorem on twisted Hilbert spaces, and proves
the theorem:
\[
  Z(T^2)\ \text{satisfies the Johnson--Lindenstrauss lemma.}
\]
The proof is by showing that \(Z(T^2)\) is log-Hilbertian and then invoking
Johnson--Naor's result that log-Hilbertian spaces satisfy the J-L lemma.

\begin{claim}
There exists a weak Hilbert space with no unconditional basis that satisfies
the Johnson--Lindenstrauss lemma.  One such space is \(Z(T^2)\).
\end{claim}

\begin{proof}[Literature implication]
By Suarez de la Fuente \cite{SdlF}, \(Z(T^2)\) is a weak Hilbert twisted
Hilbert space and has no unconditional basis.  Theorem \(\mathrm{log}\) of
\cite{SdlF} proves that \(Z(T^2)\) satisfies the J-L lemma.  Hence the class
of weak Hilbert spaces satisfying the J-L lemma strictly extends beyond the
unconditional-basis subclass mentioned in the original Johnson--Naor
discussion.
\end{proof}

\begin{scope}
This is only a partial answer.  It does not prove that every weak Hilbert
space is log-Hilbertian, \(h\)-Hilbertian with \(h(n)=o(n)\), or J-L.  It only
supplies a concrete positive example without an unconditional basis, thereby
settling that subcase/variant of the Johnson--Naor weak-Hilbert direction.
\end{scope}

\begin{thebibliography}{9}

\bibitem{JN}
W.~B. Johnson and A.~Naor,
\emph{The Johnson--Lindenstrauss lemma almost characterizes Hilbert space, but
not quite},
Discrete Comput. Geom. 43 (2010), no. 3, 542--553; arXiv:0807.1919.

\bibitem{SdlF}
J.~Suarez de la Fuente,
\emph{A space with no unconditional basis that satisfies the
Johnson--Lindenstrauss lemma},
arXiv:2501.13524.

\end{thebibliography}

\end{document}
